%%
%% introduction.tex
%% 
%% Made by Alex Nelson <pqnelson@gmail.com>
%% Login   <alex@lisp>
%% 
%% Started on  2026-03-25T09:38:49-0700
%% Last update 2026-03-25T09:38:49-0700
%% 

\section{Introduction}

\begin{node}
We implement a proof assistant which allows us to reason about metatheory.
Among our inspirations, we follow the \emph{spirit} of Hilbert's
programme and provide a proof assistant roughly corresponding to
Hilbert's ``finitary metatheory''.\footnote{More precisely, our proof assistant
is strictly not stronger than Hilbert's finitary metatheory. Tait~\cite{tait2002remarks} argues that Hilbert's finitary metatheory was precisely \PRA,
others disagree and assert that Hilbert's finitary metatheory extended
slightly beyond \PRA---for example, Zach points
out~\cite{zach2003practice} that Ackermann's consistency proofs
involved ordinals like $\omega^{\omega^{\omega}}$ which were
accepted as ``finitary'' by Hilbert and friends despite \PRA\ having a
proof ordinal strength of $\omega^{\omega}$. Since $FS_{0}$ is a conservative extension of
\PRA, whichever perspective one might take, we do not extend
\emph{beyond} \PRA.}
We use Feferman's $FS_{0}$ set theory atop Intuitionistic first-order
logic, preserving the constructive \emph{and} finitary character of
Hilbert's metatheory.

Our intention is to build a tool to facilitate \emph{exploring the
design space for proof assistants} which allows us to formally prove
results concerning our exploration. Will adding a new feature to a
proof assistant ruin the soundness of that proof assistant? Usually
such questions are proven informally, we wanted to provide
\emph{formal proofs} for them. Here we depart from Hilbert's programme
by making the informal finitary metatheory a formal proof
assistant. What is more, we provide a proof assistant which is
a \emph{metalogical framework}.
\end{node}

\begin{node}[Metalogical frameworks]
Metalogical frameworks~\cite{basin1991metalogical,basin2004reflective} allow us to state and prove metatheorems
concerning logics and proof calculi.
%A logic's syntax and proofs are built inductively from syntax and proof constructors.
Logical frameworks use these to allow us to ``reason within the logic'' and
build derivations of the object logic as objects of the metalogic. But a
metalogical framework needs \emph{more}: if the representation
preserves the inductive structure of the object logic's syntax and
proof constructors, then we can perform inductive proofs in the
metalogical framework to prove properties about the object logic.

However, the ``exact criteria'' for a metalogical framework remains
elusive. This naturally raises the concern whether ``metalogical
frameworks'' are a well-defined notion or not. The reader may rest
assured, this is well-defined because we have several examples of metalogical frameworks
like \HOL, $FS_{0}$~\cite{matthews1992experience,matthews1993theory}, and membership equational logic. How will we
recognize a metalogical framework ``in the wild''? One sufficient
criteria has been offered~\cite{basin2004reflective} as any reflective
logical framework qualifies as a metalogical framework.

Really, the critical property required for a logical framework to be a
metalogical framework is the ability to state inductive definitions
\emph{and} prove properties about inductive definitions. Moreover, we
must have the ability to do induction \emph{over} object logics. The
intuition we should have: metalogical frameworks are logical
frameworks for doing metamathematics.
\end{node}

\begin{node}[$FS_{0}$ as a metalogical framework]
Feferman's set theory $FS_{0}$ can serve as a metalogical framework.
The basic idea is that inductive definitions are terms in $FS_{0}$,
which have induction rules for reasoning about such terms. This
requires ``only'' a (three-sorted) first-order logic to state the
axioms for $FS_{0}$ and, although Feferman allowed Classical logic, we
will use an Intuitionistic logic.

This has been explored quite thoroughly in Se\'{a}n Matthews's PhD
dissertation~\cite{matthews1992metatheoretic} and related work~\cite{matthews1992experience,matthews1993theory}.
The basic idea may be found sketched out in Feferman~\cite[\S21]{feferman1989finitary}.

The major advantage favoring $FS_{0}$ as a metalogical framework is
its simplicity. \HOL\ requires building quite a bit of infrastructure
before we get to inductive data types, whereas $FS_{0}$ has them
\textit{ab initio}.

However, several major disadvantages are noted in Matthews, Smaill,
and Basin~\cite[\S5.1]{matthews1992experience}:
\begin{quote}
One obvious question to ask is `how important is the fact that we have 
implemented our system directly in Prolog?'; i.e., would $FS_{0}$ be a usable 
system if it was instead implemented using another framework such as
Isabelle or the ELF? We would argue that such an implementation would 
probably not be very useful for the following reason: $FS_{0}$ is itself intended 
as a framework theory, so if it were implemented on top of another system 
then any work done in a logic declared in it would have to work through 
two levels of encoding to get to the machine level. Further, special care was 
taken to provide an efficient evaluation mechanism for function applications, 
and this might not be available in other approaches. 
\end{quote}
Furthermore, Pfenning~\cite{pfenning2001logical} notes the
difficulties using first-order logic (plus some adequate addition like 
$FS_{0}$) as just a \emph{logical framework}. For example, determining
if a string is a well-formed formula requires quite a bit more work
than if we used \HOL.
\end{node}

\begin{node}
We are largely inspired by Lawrence Paulson's Isabelle~\cite{paulson1989foundation,paulson1990isabelle,paulson1994isabelle}.
However, whereas Isabelle is a logical framework which serves as the
``engine'' for any foundations of Mathematics, we will be interested
in a logical framework suitable for Metamathematics. The irate reader
may wonder, ``Then in what sense is Isabelle an
inspiration?''\footnote{A wag may say that our endeavour is, as
Paulson writes~\cite{paulson1990isabelle}, ``a tale of errors, not grand design.''}

Isabelle's abstractions are inspired by communication metaphors (e.g.,
scripts are \emph{documents} or \emph{articles}, comments are \emph{text}, 
\LaTeX\ is autogenerated suitable for publication, etc.). Since our
goal is to report results from exploring the design space,
communication is critically important.

Markus Wenzel's PhD dissertation~\cite{wenzel2002isabelle} offers a
principled manner of constructing ``structured'' declarative-style
proofs atop such a proof assistant. This actually spurred us on to
start this project, since Wenzel's thesis provides a specification for
Isar, and we were thinking about Mario Carneiro's
comment~\cite{carneiro2023reimplementing} that ``there is no `spec' for the Mizar language''.
What qualifies as a ``spec''? How do we communicate it? How can we
prove a proof assistant implements its specification? Trying to answer
these questions led us to this project.
\end{node}

\begin{node}[Outline of document]
We begin by implementing the abstract syntax trees for terms and
formulas for three-sorted first-order logic. Then we abstract away the
constructors for primitive notions appearing in $FS_{0}$ by providing
an \texttt{FS0} module consisting of constructors. This work
constitutes the first 30 pages or so of work.

Then we implement an LCF-style kernel for Intuitionistic natural
deduction. This includes providing the axioms for $FS_{0}$. The kernel
alone (i.e., the \texttt{Thm} module) consists of about 320 lines of
code. Since the kernel is minimal, we provide a larger
\textit{r\'{e}pertoire} of theorems and derived inference
rules. Altogether, this constitutes the next 54 pages or so.

We then implement the LCF-style tactics for proving theorems. 
\end{node}

\begin{node}
The appendices contains code used throughout the program, but are a
bit of a ``distraction'' from the reader interested exclusively in
proof assistants.
\end{node}